\documentclass[12pt]{article}
\usepackage[utf8]{inputenc}
\usepackage[T1]{fontenc}
\usepackage{hyperref}
\title{Zapotec, Lapagu\'ia-Guivini Language Resources for K-12 Curriculum Development}
\author{Compiled by Claude Sonnet 5.5 (claude-sonnet-5-5) for Andr\'as Kornai}
\date{}
\begin{document}
\maketitle

\section{Introduction}
Lapagu\'ia-Guivini Zapotec (ISO 639-3 ztl, Glottocode lapa1249) is an Otomanguean language of southern Oaxaca, Mexico. Glottolog places it in Central Core Zapotec. Speaker counts in the sources I found disagree: Wikipedia gives about 4,200 (citing 1983), Global Recordings Network gives 6,840, and Joshua Project gives about 7,700. Very little educational material turned up in my searches. This report lists only what the searches returned.

\section{Primary Educational Resources}
No primers, textbooks or school programmes for this language were found in the searches. Wikitongues lists the language as unwritten, so a standard orthography appears not to exist.

\section{Linguistic and Reference Materials}
\begin{itemize}
\item Wikipedia overview (dialects Lapagu\'ia and Guivini): \url{https://en.wikipedia.org/wiki/Lapagu%C3%ADa_Zapotec}
\item Glottolog entry: \url{https://glottolog.org/resource/languoid/id/lapa1249.bigmap.html}
\item Wikitongues language page: \url{https://wikitongues.org/languages/ztl/}
\end{itemize}
No dictionary or grammar was found.

\section{Religious and Community Resources}
\begin{itemize}
\item Scripture Earth listing (New Testament audio): \url{https://www.scriptureearth.org/00eng.php?iso=ztl}
\item Joshua Project people-group profile: \url{https://www.joshuaproject.net/people_groups/14745/MX} (a missions organisation, so its framing is religious)
\end{itemize}

\section{Digital and Audio Resources}
\begin{itemize}
\item Global Recordings Network audio (Bible stories): \url{https://globalrecordings.net/en/language/ztl}
\end{itemize}

\section{Curriculum Development Context}
Little documentation appears to exist. A curriculum would need new work first: community-led orthography development, a lexicon and a recorded corpus. Audio recordings and Scripture audio could serve as starting texts, with community speakers involved.
\end{document}